%! Author = lili
%! Date = 5/23/26


\section{Types of Knowledge}
Different levels of knowledge that are central to understanding how people learn and
reason:

\defin{declarative-knowledge}
\defin{structural-knowledge}
\defin{procedural-knowledge}

\begin{bsp}
    \textbf{Examples for \autoref{def:declarative-knowledge}, \autoref{def:structural-knowledge} and
    \autoref{procedural-knowledge}} \\
    \begin{itemize}
        \item \Gls{declarative-knowledge}: warm air rises
        \item \Gls{structural-knowledge}: heat $\rightarrow$ low density $\rightarrow$ buoyancy
        \item \Gls{procedural-knowledge}: performing a convection calculation
    \end{itemize}
\end{bsp}

\subsection*{What is Structural Knowledge?}

\begin{center}
    `` \textit{In order to know how, you must know why. Structural knowledge provides the
    conceptual bases for why; it describes how the declarative knowledge is
    interconnected} ''
\end{center}
(Jonassen, Beissner, and Yacci 1993)

According to Jonassen, Beissner, and Yacci (1993), the meaning of a concept is implicit in its pattern of relationships to other concepts.
Therefore, structural knowledge is closely related to conceptual knowledge, internal connectedness, and integrative understanding.
Tennyson and Cocchiarella (1986) define conceptual knowledge as the integrated storage of meaningful dimensions within a specific knowledge domain.
Structural knowledge also plays an important role in learning because it mediates the acquisition of procedural
knowledge by connecting facts and concepts to actions and procedures.

\subsection*{Why Not Just First-Order Logic?}
\Gls{fol} quantifies over individual objects only, not over predicates or sets.
This creates fundamental expressive limits:
\begin{itemize}
    \item Cannot quantify over predicates or sets. Anything about sets must be
    rephrased in terms of predicates, which is often not possible
    (Quora contributors, 2024).

    \item Cannot define mathematical induction in a single finite sentence.
    Induction requires quantifying over predicates --- a second-order principle
    (Wikipedia contributors, 2026b; Logic and Proof contributors, 2024).

    \item Cannot express ``finitely many'' or characterise finiteness/countability:
    by the Löwenheim--Skolem theorem, any consistent FOL theory with an infinite
    model also has a countable model (Wikipedia contributors, 2026b).

    \item No natural grouping of properties around a concept, no built-in defaults,
    and no context (Steels, 1979).
\end{itemize}

\begin{lem}
    \textbf{FOL requires an axiom schema}
    FOL cannot write a single axiom
    $\forall P.(P (0) \land \forall k. (P (k) \rightarrow P (k+1))) \rightarrow \forall n. P (n)$
    because $\forall P$
    quantifies over predicates. Instead, FOL needs one axiom per predicate (wiki 2026b).
\end{lem}

\begin{theo}
    \textbf{Constructing a knowledge representation} \\
    ``\textit{What we do in constructing a knowledge representation language is to make an
    abstraction over a class of representational structures and introduce a syntactic
    mechanism to represent that abstraction.}'' \\
    (Steels 1979)
\end{theo}


\section{Frames and Slots} % 9/44

\defin{frame}


\section{Object-Oriented Analogy: Classes and Instances} % 18/44


\section{Inheritance Hierarchies} % 23/44


\section{Scripts and Semantic Networks} % 27/44


\section{Default Reasoning and Conditionals} % 30/44


\section{Problems with Schema Systems} % 33/44


\section{Frames in Practice: JSON-LD} % 35/44

